\clearpage
\oldpage[IV]{345}
\phantomsection
\addcontentsline{toc}{section}{Errata and Addenda (List 3)}
\section*{Errata and Addenda}
\begin{center}\large (List 3)\end{center}

\subsection*{A) Typographical errors}

\paragraph{(\textbf{II}, 3.3.2)}
On line 1 of p.~58, replace $\sh{S}$ by $\sh{M}$.  On line 6 of p.~58,
replace $S(n)$ by $M(n)$.

\paragraph{(\textbf{II}, 4.1.1)}
On line 8 of p.~71, replace $\sh{O}$-${}_P$Module by
$\sh{O}_P$-Module.

\paragraph{(\textbf{II}, 5.4.5)}
On line 6 of p.~102, replace $u$ by $f\circ u$.

\paragraph{(\textbf{II}, 5.4.8)}
Before the fourth line from the bottom of p.~102, replace
$f\times 1_{S'}$ by $f_{(Y')}$ in the diagram.  On the fourth line from
the bottom of p.~102, replace $f\times 1_{S'}$ by $f_{(Y')}$.

\paragraph{(\textbf{II}, 6.4.4)}
On the first line at the top of p.~122, replace ``normal'' by
``integrally closed.''

\paragraph{(\textbf{II}, 6.4.7)}
On the fourteenth line from the bottom of p.~122, replace ``dimension''
by ``rank.''

\paragraph{(\textbf{II}, 6.5.2)}
On line 5 of p.~127, replace $(\sh{B}|U_\lambda)$ by $(B|U_\lambda)$.
On line 15 of p.~127, delete the first occurrence of
$(N_{\sh{B}/\sh{A}}\omega_{\lambda_\mu})$.

\paragraph{($\mathbf{0}_{\mathrm{III}}$, 8.1.11)}
On the eighth line from the bottom of p.~8, replace $h_X(u)$ by
$h'_X(u)$.

\paragraph{($\mathbf{0}_{\mathrm{III}}$, 9.1.2)}
On the tenth line from the bottom of p.~12, replace $\mathfrak{F}$ by $X$.

\paragraph{($\mathbf{0}_{\mathrm{III}}$, 10.2.1)}
On the sixth line from the bottom of p.~18, replace ``in addition if''
by ``if in addition.''

\paragraph{(\textbf{III}, 4.4.12)}
On line 10 of p.~138, replace ``chap.~V'' by ``chap.~IV (8.11.2).''

\paragraph{(\textbf{III}, 5.2.5)}
On line 12 of p.~153, replace (\textbf{I}, 5.5.4) by (\textbf{II}, 5.5.4).

\paragraph{(\textbf{III}, 7.4.7)}
On line 15 of p.~57, replace ``every $A$-module'' by
``every $A$-module $M$.''

\paragraph{(\textbf{III}, 7.8.5)}
On the eighteenth line from the bottom of p.~74, replace $F$ by $\sh{F}$.

\paragraph{(\textbf{III}, 7.8.9)}
On the ninth line from the bottom of p.~75, replace $M$ by $\sh{M}$.

\paragraph{(\textbf{III}, 7.9.11)}
On the seventeenth line from the bottom of p.~79, replace chap.~V by
chap.~VI.

\paragraph{($\mathbf{0}_{\mathrm{IV}}$, 16.2.2.2)}
On line 13 of p.~26, replace $M'\cap\mathfrak q^nM'$ by
$M'\cap\mathfrak q^nM$.

\paragraph{($\mathbf{0}_{\mathrm{IV}}$, 17.3.7)}
On the sixteenth line from the bottom of p.~50, replace (17.1.10) by
(17.3.3).

\paragraph{($\mathbf{0}_{\mathrm{IV}}$, 18.1.3)}
In diagram (18.1.3.1), $p_2$ must be beside the arrow $G\to F$.

\paragraph{($\mathbf{0}_{\mathrm{IV}}$, 19.8.7)}
On the seventh, eighth, and ninth lines from the bottom of p.~111,
replace (21.5.5) by (21.5.3).

\paragraph{($\mathbf{0}_{\mathrm{IV}}$, 19.8.8)}
On the second line from the bottom of p.~112, replace $A$-module by
$B$-module.

\paragraph{($\mathbf{0}_{\mathrm{IV}}$, 20.5.4)}
On the fifth line from the bottom of p.~130, replace
$\Omega_{B'/A'}$ by $\Omega^1_{B'/A'}$.

\paragraph{($\mathbf{0}_{\mathrm{IV}}$, 20.7.8)}
On the ninth line from the bottom of p.~149, replace (18.4.3) by
(18.5.4).  On the fourth line from the bottom, replace (18.4.1) by
(18.5.1).

\paragraph{($\mathbf{0}_{\mathrm{IV}}$, 21.5.3)}
On the last line of p.~166, replace the feminine form of ``ramified'' by
the masculine form.

\paragraph{($\mathbf{0}_{\mathrm{IV}}$, 21.7.1)}
On line 4 of p.~170, replace (20.6.15.1) by (20.6.17.1).

\paragraph{($\mathbf{0}_{\mathrm{IV}}$, 22.1.1)}
On line 12 of p.~183, replace (20.6.20) by (20.6.10).

\oldpage[IV]{346}
\paragraph{($\mathbf{0}_{\mathrm{IV}}$, 22.3.3)}
On the thirteenth line from the bottom of p.~192, replace
$\widehat B_{\mathfrak p}$ by $\widehat B_{\mathfrak q}$.

\paragraph{($\mathbf{0}_{\mathrm{IV}}$, 22.7.3)}
On the twenty-first line from the bottom of p.~212, replace
$\operatorname{Ker}(i)$ by $\operatorname{Ker}(i_0)$.  On the twelfth
line from the bottom of p.~212, replace
$(\mathfrak p/\mathfrak p^2)\otimes_L L$ by
$(\mathfrak p/\mathfrak p^2)\otimes_A L$.

\paragraph{($\mathbf{0}_{\mathrm{IV}}$, 23.1.3)}
On line 4 of p.~215, replace $x'^q\in K'$ by $x'\in K'$.  On line 9,
replace $x'^q\in A'$ by $x'^q\in A$.

\paragraph{({\texorpdfstring{\protect\hyperref[IV.2.1.1]{\textbf{IV}, 2.1.1}}{IV, 2.1.1}})}
On the fourth line from the bottom of p.~5, replace (\textbf{III}, 1.4.15.1) by
(\textbf{III}, 1.4.15.5).

\paragraph{({\texorpdfstring{\protect\hyperref[IV.2.4.3]{\textbf{IV}, 2.4.3}}{IV, 2.4.3}})}
On line 12 of p.~20, replace
$(X_{\mathrm{red}}\times_{Y_{\mathrm{red}}}Y_{\mathrm{red}})_{\mathrm{red}}$
by
$(X_{\mathrm{red}}\times_{Y_{\mathrm{red}}}Y'_{\mathrm{red}})_{\mathrm{red}}$.

\paragraph{({\texorpdfstring{\protect\hyperref[IV.2.6.1]{\textbf{IV}, 2.6.1}}{IV, 2.6.1}})}
On line 10 of p.~27, replace $X'\to X_{(S')}$ by $X'=X_{(S')}$.

\paragraph{({\texorpdfstring{\protect\hyperref[IV.4.2.1]{\textbf{IV}, 4.2.1}}{IV, 4.2.1}})}
On the twentieth line from the bottom of p.~55, replace
$\operatorname{tr.deg}_kE$ by $\operatorname{tr.deg}_LE$ and
$\operatorname{tr.deg}_kL(\mathbf{t})$ by $\operatorname{tr.deg}_LL(\mathbf{t})$; on the
nineteenth line from the bottom, interchange ``first'' and ``second.''

\paragraph{({\texorpdfstring{\protect\hyperref[IV.4.4.2]{\textbf{IV}, 4.4.2}}{IV, 4.4.2}})}
On the last line of p.~59, replace $X$ by $X'$.

\paragraph{({\texorpdfstring{\protect\hyperref[IV.4.7.3]{\textbf{IV}, 4.7.3}}{IV, 4.7.3}})}
On the tenth line from the bottom of p.~75, replace ``prescheme'' by
``$k$-prescheme.''

\paragraph{({\texorpdfstring{\protect\hyperref[IV.5.6.1]{\textbf{IV}, 5.6.1}}{IV, 5.6.1}})}
On line 16 of p.~98, replace the subscript $\boldsymbol{k}(\mathfrak p($ by
$\boldsymbol{k}(\mathfrak p)$.

\paragraph{({\texorpdfstring{\protect\hyperref[IV.5.6.5]{\textbf{IV}, 5.6.5}}{IV, 5.6.5}})}
On line 20 of p.~99, delete the second closing parenthesis in $\boldsymbol{k}(\xi))$.

\paragraph{({\texorpdfstring{\protect\hyperref[IV.5.6.11.1]{\textbf{IV}, 5.6.11.1}}{IV, 5.6.11.1}})}
On line 16 of p.~102, replace $K$ by $\mathfrak K$.

\paragraph{({\texorpdfstring{\protect\hyperref[IV.5.11.2]{\textbf{IV}, 5.11.2}}{IV, 5.11.2}})}
In the footnote on p.~123, replace ``5.11.2)'' by ``(5.11.2).''

\paragraph{({\texorpdfstring{\protect\hyperref[IV.6.1.5]{\textbf{IV}, 6.1.5}}{IV, 6.1.5}})}
On the second line from the bottom of p.~136, replace (5.5.1.2) by
(\textbf{0}, 16.3.9.2).

\paragraph{({\texorpdfstring{\protect\hyperref[IV.6.3.1]{\textbf{IV}, 6.3.1}}{IV, 6.3.1}})}
On line 9 of p.~139, replace $y\in\mathfrak m$ by $y\in\mathfrak n$.

\paragraph{({\texorpdfstring{\protect\hyperref[IV.6.7.2]{\textbf{IV}, 6.7.2}}{IV, 6.7.2}})}
On line 18 of p.~146, replace $\sh{O}'_{x'}$ by $\sh{O}_{x'}$.

\paragraph{({\texorpdfstring{\protect\hyperref[IV.6.10.1]{\textbf{IV}, 6.10.1}}{IV, 6.10.1}})}
On line 13 of p.~155, replace
$\sh{F}\otimes_{\sh{O}_Y}\sh{O}_Y$ by
$\sh{F}\otimes_{\sh{O}_X}\sh{O}_Y$.

\paragraph{({\texorpdfstring{\protect\hyperref[IV.6.14.1]{\textbf{IV}, 6.14.1}}{IV, 6.14.1}})}
On line 21 of p.~170, replace ``no.~3, prop.~12'' by
``no.~2, prop.~8.''

\paragraph{({\texorpdfstring{\protect\hyperref[IV.6.15.2]{\textbf{IV}, 6.15.2}}{IV, 6.15.2}})}
On line 13 of p.~177, replace $t(t\setmin v)$ by $t^2(t\setmin v)^2$.

\paragraph{({\texorpdfstring{\protect\hyperref[IV.6.15.7]{\textbf{IV}, 6.15.7}}{IV, 6.15.7}})}
On the sixth line from the bottom of p.~179, replace
$X\otimes_Ak'$ by $X\otimes_kk'$.

\paragraph{({\texorpdfstring{\protect\hyperref[IV.7.5.4]{\textbf{IV}, 7.5.4}}{IV, 7.5.4}})}
On the fourteenth line from the bottom of p.~206, replace
$\operatorname{Spec}(\widehat B\otimes k)$ by
$\operatorname{Spec}(\widehat B\otimes_{\widehat A}k)$.

\paragraph{({\texorpdfstring{\protect\hyperref[IV.7.8.4]{\textbf{IV}, 7.8.4}}{IV, 7.8.4}})}
On the tenth line from the bottom of p.~216, move the opening parenthesis
so that the phrase reads ``satisfying ({\hyperref[IV.7.8.2]{7.8.2}}, (ii) and (iii)).''

\paragraph{(\textbf{IV}, Part 2, Bibliography)}
On line 6 of p.~224, replace \emph{Amer. J. of Math.} by
\emph{Annals of Math.}

\paragraph{({\texorpdfstring{\protect\hyperref[IV.8.5.8]{\textbf{IV}, 8.5.8}}{IV, 8.5.8}})}
On line 6 of p.~24, replace $\operatorname{Ker}g_\mu$ at the end of the
line by $\operatorname{Ker}g_\lambda$.

\paragraph{({\texorpdfstring{\protect\hyperref[IV.8.10.1]{\textbf{IV}, 8.10.1}}{IV, 8.10.1}})}
On the thirteenth line from the bottom of p.~36, replace (2.3.11) by
{\hyperref[IV.8.2.11]{(8.2.11)}}.

\paragraph{({\texorpdfstring{\protect\hyperref[IV.10.2.1]{\textbf{IV}, 10.2.1}}{IV, 10.2.1}})}
On the fourteenth line from the bottom of p.~98, replace $y$ by $Y$.

\paragraph{({\texorpdfstring{\protect\hyperref[IV.10.4.11]{\textbf{IV}, 10.4.11}}{IV, 10.4.11}})}
On the tenth line from the bottom of p.~104, replace (17.9.7) by
{\hyperref[IV.17.9]{(17.9.6)}}.

\paragraph{({\texorpdfstring{\protect\hyperref[IV.11.2.9]{\textbf{IV}, 11.2.9}}{IV, 11.2.9}})}
On the third line from the bottom of p.~126, replace $A_\lambda$ by
$A_\mu$.

\paragraph{({\texorpdfstring{\protect\hyperref[IV.11.3.4]{\textbf{IV}, 11.3.4}}{IV, 11.3.4}})}
On line 16 of p.~134, replace $F$ by $\sh{F}$.

\paragraph{({\texorpdfstring{\protect\hyperref[IV.11.3.10]{\textbf{IV}, 11.3.10}}{IV, 11.3.10}})}
On the ninth line from the bottom of p.~139, replace (11.3.7) by
{\hyperref[IV.11.3.10]{(11.3.10)}}.

\paragraph{({\texorpdfstring{\protect\hyperref[IV.11.7.5]{\textbf{IV}, 11.7.5}}{IV, 11.7.5}})}
On line 11 of p.~159, replace (6.5.11) by (6.15.11).

\paragraph{({\texorpdfstring{\protect\hyperref[IV.12.1.1]{\textbf{IV}, 12.1.1}}{IV, 12.1.1}})}
On line 13 of p.~174, add at the end: ``for $k\geq 1$.''

\paragraph{({\texorpdfstring{\protect\hyperref[IV.14.2.1]{\textbf{IV}, 14.2.1}}{IV, 14.2.1}})}
On line 11 of p.~202, replace (6.1.2) by (6.1.1).

\oldpage[IV]{347}
\paragraph{({\texorpdfstring{\protect\hyperref[IV.14.3.8]{\textbf{IV}, 14.3.8}}{IV, 14.3.8}})}
On the tenth line from the bottom of p.~206, replace $(f^{-1}y)$ by
$f^{-1}(y)$.

\paragraph{({\texorpdfstring{\protect\hyperref[IV.14.5.7]{\textbf{IV}, 14.5.7}}{IV, 14.5.7}})}
On line 16 of p.~218, replace (14.5.5) by (14.5.6).

\paragraph{({\texorpdfstring{\protect\hyperref[IV.14.5.11]{\textbf{IV}, 14.5.11}}{IV, 14.5.11}})}
In the diagram at the top of p.~223, every horizontal arrow must point
from right to left.

\paragraph{({\texorpdfstring{\protect\hyperref[IV.15.4.2]{\textbf{IV}, 15.4.2}}{IV, 15.4.2}})}
On line 23 of p.~230, replace (6.1.4) by (6.1.5).

\paragraph{({\texorpdfstring{\protect\hyperref[IV.15.5.2]{\textbf{IV}, 15.5.2}}{IV, 15.5.2}})}
On line 16 of p.~233, replace $f^{-1}(y)$ by $f^{-1}(y')$.

\paragraph{({\texorpdfstring{\protect\hyperref[IV.15.6.10]{\textbf{IV}, 15.6.10}}{IV, 15.6.10}})}
On line 6 of p.~242, replace ${}^{-1}(y')$ by $f^{-1}(y')$.  On line
11, replace ``$\alpha$) and $\beta$)'' by ``$\alpha$) or $\beta$).''

\subsection*{B) Textual modifications}

\paragraph{($\mathbf{Err}_{\mathrm{IV}}$, 1)}
In ($\mathbf{0}_{\mathrm{III}}$, 11.1.6), the statement (which reproduces a
statement of (G, I, 4.4.1)) is false when $q(n)$ is an arbitrary integer;
it is correct only when $q(n)=0$ or $q(n)=n$.  Indeed, in the first case
the kernel of $d_r^{p,0}$ is $E_r^{p0}$, since the image of
$d_r^{p,0}$ lies in $E_r^{p+r,-r+1}=0$ by hypothesis; likewise the
image of $d_r^{p-r,r-1}$ is zero since $E_r^{p-r,r-1}=0$.  Hence
$E_{r+1}^{p,0}=E_r^{p,0}$.  The case $q(n)=n$ for every $n$ is handled
in the same way.

\paragraph{($\mathbf{Err}_{\mathrm{IV}}$, 2)}
In (\textbf{III}, 3.4), delete lines 2 and 3 of p.~119.

\oldpage[IV]{348}
\paragraph{(Err\textsubscript{IV}, 3).}
In (\textbf{III}, 4.4.9), on the fifteenth line from the bottom of p.~137,
replace ``finite'' by ``finite as sets.''

\paragraph{($\mathbf{Err}_{\mathrm{IV}}$, 4)}
In (\textbf{III}, 7.8.10), on line 9 of p.~76, replace ``simple'' by ``smooth.''

\paragraph{($\mathbf{Err}_{\mathrm{IV}}$, 5)}
In ($\mathbf{0}_{\mathrm{IV}}$, 16.1.4), on the eighth line from the bottom of
p.~23, replace ``To say that $A$ is catenary'' by ``To say that an
\emph{integral} ring $A$ is catenary.''

\paragraph{($\mathbf{Err}_{\mathrm{IV}}$, 6)}
In ($\mathbf{0}_{\mathrm{IV}}$, 16.1.9), the hypothesis that $A$ is Noetherian
may be deleted.  Indeed, for every $u\in B_M$, $A[u]$ is an
$A$-subalgebra of $C$, so $u$ is integral over $A$ (Bourbaki,
\emph{Alg. comm.}, chap.~V, \S1, no.~1, corollary to prop.~1), and
hence also over $A_M$; one concludes again by {\hyperref[section:IV.16]{(16.1.5)}}.

\paragraph{($\mathbf{Err}_{\mathrm{IV}}$, 7)}
In ($\mathbf{0}_{\mathrm{IV}}$, 16.5.12), note that relation (16.5.11.1) is
equivalent to the same relation after replacing $A$ by the quotient rings
$A/\mathfrak r$, where $\mathfrak r$ ranges over the minimal prime ideals
of $A$.  Thus in the proof of {\hyperref[section:IV.16]{(16.5.12)}} one may reduce to the case in
which $A$ is integral and apply (16.1.4.2).

\paragraph{($\mathbf{Err}_{\mathrm{IV}}$, 8)}
Delete ($\mathbf{0}_{\mathrm{IV}}$, 19.3.12), whose proof is erroneous (the sets
$W_\lambda$, $U_\lambda$, and $V_\lambda$ are not linear varieties).

\paragraph{($\mathbf{Err}_{\mathrm{IV}}$, 9)}
In ($\mathbf{0}_{\mathrm{IV}}$, 19.8.3), on the sixth and eighth lines from the
bottom of p.~110, replace ``homomorphism'' by ``local homomorphism'';
make the same replacement in ($\mathbf{0}_{\mathrm{IV}}$, 19.8.4), line 1 of
p.~111.

\paragraph{($\mathbf{Err}_{\mathrm{IV}}$, 10)}
In ($\mathbf{0}_{\mathrm{IV}}$, 20.1.2), on the second line from the bottom of
p.~117, replace $a\in A$ by $b\in B$ and $a\mathrm{D}$ by $b\mathrm{D}$; on the last line,
replace $A$-module by $B$-module.

\paragraph{($\mathbf{Err}_{\mathrm{IV}}$, 11)}
In ($\mathbf{0}_{\mathrm{IV}}$, 20.2.1), on line 8 of p.~119, replace
$A$-modules by $B$-modules.  On lines 21 and 22, replace
``di-homomorphisms of modules (relative to $u$)'' by ``homomorphism of
$B$-modules.''

\paragraph{($\mathbf{Err}_{\mathrm{IV}}$, 12)}
In ($\mathbf{0}_{\mathrm{IV}}$, 20.4.6.1), replace $p_1\setmin p_2$ by $p_2\setmin p_1$; on
the tenth line from the bottom of p.~125, replace
$x\otimes1\setmin 1\otimes x$ by $1\otimes x\setmin x\otimes1$.

\paragraph{(Err\textsubscript{IV}, 13).}
In ($\mathbf{0}_{\mathrm{IV}}$, 20.4.8), on line 14 of p.~126, replace
$p_2=p_1\setmin d_{B/A}$ by $p_1=p_2\setmin d_{B/A}$.

\paragraph{($\mathbf{Err}_{\mathrm{IV}}$, 14)}
In ($\mathbf{0}_{\mathrm{IV}}$, 20.4.13), replace the sixth through eleventh
lines from the bottom of p.~127 by the following text:

\medskip
\noindent
(iv) Let $A$ and $B$ be integral rings such that $A\subset B$, with
$B$ integral over $A$, and suppose that the fraction field of $B$ is a
separable extension of that of $A$.  Then the $B$-module
$\Omega^1_{B/A}$ is a torsion module.  Indeed, for every $x\in B$, let
$f$ be the minimal polynomial of $x$ over the fraction field $K$ of
$A$; there is a nonzero $a\in A$ such that $af[T]\in A[T]$.  Since $x$
is separable over $K$, $f'(x)\neq0$, while the relation $af(x)=0$
gives $af'(x)d_{B/A}x=0$; the assertion follows from {\hyperref[IV.20]{(20.4.7)}}.

\paragraph{($\mathbf{Err}_{\mathrm{IV}}$, 15)}
In ($\mathbf{0}_{\mathrm{IV}}$, 22.2.7), on line 5 of p.~189, replace
``formally smooth'' by ``formally smooth for the preadic topologies.''

\paragraph{($\mathbf{Err}_{\mathrm{IV}}$, 16)}
In the Summary of Chapter IV, on the fourth line from the bottom of
p.~222, add ``and strictly local rings.''  Replace the last three lines
by:
\begin{quote}
\S19. Regular immersions and normal flatness.\\
\S20. Meromorphic functions and pseudo-morphisms.\\
\S21. Divisors.
\end{quote}
On line 16 of p.~223, replace ``of finite presentation'' by ``locally of
finite presentation''; on line 19, replace ``as'' by ``as being.''

\paragraph{($\mathbf{Err}_{\mathrm{IV}}$, 17)}
In (\textbf{IV}, 3.1.10), add after line 16 of p.~39:

\begin{remark}[3.1.10.2]
\label{IV.3.1.10.2}
Let $f:X\to Y$ be a quasi-compact and quasi-separated morphism, and let
$\sh{F}$ be a quasi-coherent $\sh{O}_X$-Module, so that
$f_*(\sh{F})$ is a quasi-coherent $\sh{O}_Y$-Module
{\hyperref[IV.1.7.4]{(1.7.4)}}. Then we have
\begin{equation}
 \operatorname{Ass}(f_*(\sh{F}))\subset f(X)
 \label{IV.3.1.10.3}\tag{3.1.10.3}
\end{equation}
and, if $X$ is locally Noetherian,
\begin{equation}
 \operatorname{Ass}(f_*(\sh{F}))\subset f(\operatorname{Ass}(\sh{F})).
 \label{IV.3.1.10.4}\tag{3.1.10.4}
\end{equation}

Indeed, let $y\in\operatorname{Ass}(f_*(\sh{F}))$; by virtue of
{\hyperref[IV.2.3.1]{(2.3.1)}} and of the fact that the morphism
$\operatorname{Spec}(\sh{O}_{Y,y})\to Y$ is flat,
$(f_*(\sh{F}))_y$ is equal to $(f'_*(\sh{F}'))_y$, where
$f':X\times_Y\operatorname{Spec}(\sh{O}_{Y,y})\to\operatorname{Spec}(\sh{O}_{Y,y})$
is deduced from $f$ by base change and
$\sh{F}'=\sh{F}\otimes_{\sh{O}_Y}\sh{O}_{Y,y}$. We are therefore reduced to
proving (3.1.10.3) and (3.1.10.4) with $Y$ replaced by
$\operatorname{Spec}(\sh{O}_{Y,y})$, $f$ by $f'$ and $\sh{F}$ by
$\sh{F}'$, taking into account (\textbf{I}, 3.6.5) and the definition of
$\operatorname{Ass}(\sh{F})$.
By definition, there is a nonzero element
$t_y\in(f_*(\sh{F}))_y$ whose annihilator is $\mathfrak m_y$, and since
$(f_*(\sh{F}))_y$ is identified with $\Gamma(Y,f_*(\sh{F}))$ because $Y$ is
local, $t_y$ is also identified with a section $s\in\Gamma(X,\sh{F})$. The
annihilator Ideal of $s$ in $\sh{O}_X$ therefore contains $\mathfrak m_y\sh{O}_X$,
and consequently ($\mathbf{0}_{\mathrm I}$, 1.7.4), the support of $s$ (the set of
$x\in X$ such that $s_x\neq0$) is contained in
$f^{-1}(y)$; if the latter were empty, we would have $s=0$, which is
absurd. When moreover $X$ is locally Noetherian, the coherent
sub-$\sh{O}_X$-Module $\sh{O}_Xs$ of $\sh{F}$ is not zero,
hence $\operatorname{Ass}(\sh{O}_Xs)\neq\varnothing$
{\hyperref[IV.3.1.5]{(3.1.5)}}; but
$\operatorname{Ass}(\sh{O}_Xs)\subset\operatorname{Supp}(\sh{O}_Xs)\subset f^{-1}(y)$ and on the other hand
$\operatorname{Ass}(\sh{O}_Xs)\subset\operatorname{Ass}(\sh{F})$
{\hyperref[IV.3.1.7]{(3.1.7)}}, hence
$y\in f(\operatorname{Ass}(\sh{F}))$.
\end{remark}

\long\def\EGAIVErrataThirtyTwoThroughFiftyThree{%
\paragraph{(Err\textsubscript{IV}, 32).}
In (\textbf{IV}, 11.3.14), on line 7 of p.~141, insert ``(i)'' before ``Let.''
After line 9 of p.~141, add:

\smallskip
\noindent
\textup{(ii)} Conversely, let $f:X\to Y$ be a flat morphism.  If $X$
is unibranch (respectively geometrically unibranch) at a point $x$, then
$Y$ is unibranch (respectively geometrically unibranch) at the point
$y=f(x)$.

At the beginning of line 10 of p.~141, insert ``(i).''  After line 20
of p.~141, add:

\smallskip
\noindent
\textup{(ii)} Put $Y_1=Y_{\mathrm{red}}$ and
$X_1=X\times_Y Y_1$.  The underlying space of $X_1$ is that of $X$,
and the morphism $f_1:X_1\to Y_1$ obtained from $f$ by base change is
flat.  We may therefore restrict to the case
$Y=\operatorname{Spec}(A)$, with $A$ a reduced local ring and $y$ the
closed point of $Y$.  Put $B=\sh{O}_{X,x}$.  By hypothesis
$B_{\mathrm{red}}$ is unibranch, so $\operatorname{Spec}(B)$ has only
one maximal point; flatness {\hyperref[IV.2.3.4]{(2.3.4)}} then gives the same conclusion for
\oldpage[IV]{356}
$\operatorname{Spec}(A)$.  Since $A$ is reduced, it is therefore
integral.

Let $K$ be the fraction field of $A$, and let $A'$ be its integral
closure.  By flatness, $B\otimes_A A'$ identifies with a subring of
$B\otimes_A K$ containing $B$.  Since the nonzero elements of $A$ are
$B$-regular by flatness, $B\otimes_A K$ identifies with a subring of the
total ring of fractions $R$ of $B$.  To show that $A$ is unibranch, it
is enough to show that $A'$ cannot have two distinct maximal ideals
(necessarily lying over $y$; Bourbaki, \emph{Alg. comm.}, chap.~V,
\S2, no.~1, prop.~1).  If it did, then
$B'=B\otimes_A A'$ would have two distinct prime ideals over $x$
(\textbf{I}, 3.4.7), and we need only show that this is impossible.

Let $\mathfrak N$ be the nilradical of $B$.  If $S$ is the set of
regular elements of $B$, then $R=S^{-1}B$, and
$S^{-1}\mathfrak N=R\mathfrak N$ is the nilradical of $R$.  Hence
\[
 R/R\mathfrak N=S^{-1}B/S^{-1}\mathfrak N
 =S^{-1}(B/\mathfrak N)
\]
is contained in the fraction field $L$ of the integral local ring
$B/\mathfrak N=B_{\mathrm{red}}$.  The same is therefore true of
\[
 B'_{\mathrm{red}}=B'/(B'\cap R\mathfrak N),
\]
which is integral over $B_{\mathrm{red}}$.  Since
$B'_{\mathrm{red}}$ is contained in the integral closure of
$B_{\mathrm{red}}$, that ring is local by hypothesis; hence
$B'_{\mathrm{red}}$, and therefore $B'$, is \emph{local}.  This proves the
assertion.

If $B$ is moreover geometrically unibranch, the residue field of
$B'_{\mathrm{red}}$, being contained in the residue field of the
integral closure of $B_{\mathrm{red}}$, is a radicial extension of
$\boldsymbol{k}(x)$.  But the residue field of $B'_{\mathrm{red}}$ is also that of
$B'=B\otimes_A A'$, hence is
$\boldsymbol{k}(x)\otimes_{\kres(y)}k'$, where $k'$ is the residue field of $A'$.  It
follows from {\hyperref[IV.2.6.1]{(2.6.1)}} that $k'/\boldsymbol{k}(y)$ is radicial, so $A$ is geometrically
unibranch as well.

\paragraph{(Err\textsubscript{\textmd{IV}}, 33).}
In ({\hyperref[IV.12.1.2]{\textbf{IV}, 12.1.2}}), the example given
in (i) is not correct, the morphism considered being open. We
construct an example of a flat morphism $f:X\to Y$ whose restriction
to an irreducible component of $X$ is not an open morphism in
the following way. Let $Z=\operatorname{Spec}(\mathbb C[U,V])$ be the
affine plane over the field of complex numbers, and consider the
scheme $Y$ obtained by gluing together two distinct closed points $a$, $b$ of
$Z$; we have $Y=\operatorname{Spec}(B)$, where $B$ is the subring of
$\mathbb C[U,V]$ formed by the polynomials taking the same value at the points
$a$ and $b$ (the set of closed points of $Z$ being identified with
$\mathbb C^2$); it is immediate that $B$ is a sub-$\mathbb C$-vector
space of $\mathbb C[U,V]$ of codimension 1, hence $\mathbb C[U,V]$
is a $B$-module of finite type, and consequently the morphism $Z\to Y$ is
finite, the inverse image of the point $c$ of $Y$ obtained by gluing
$a$ and $b$ together being the set $\{a,b\}$.
Consider then the scheme $X$ obtained by gluing two copies of
$Z$, say $Z_1$, $Z_2$, in such a way that if $a_i$, $b_i$ ($i=1,2$) are
the points of $Z_i$ corresponding to $a$ and $b$, $a_1$ is glued to
$b_2$ and $b_1$ to $a_2$ (cf. ({\hyperref[IV.6.5.5]{6.5.5}},
(ii))). One sees easily \emph{(loc. cit.)} that the morphism $X\to Y$
thus obtained is flat and finite (and even \emph{\'etale}) and that $Z_1$ and $Z_2$
are identified with the two irreducible components of $X$. However, if
$D$ is a line in $Z_1$ passing through $a_1$ but not through $b_1$,
the image $V$ in $Y$ of the open set $Z_1\setmin D$ of $Z_1$ contains the image of
$b_1$, that is to say the point $c$, but does not contain the image of the
generic point of $D$, hence $V$ is not a neighbourhood of $c$.

\paragraph{(Err\textsubscript{\textmd{IV}}, 34).}
In the title of (\textbf{IV}, 13.2), on the fifth line from the bottom of p.~190,
replace ``equidimensional'' by ``locally equidimensional.''  Likewise,
in (\textbf{IV}, 13.2.2), on line 4 of p.~191, replace both occurrences of
``equidimensional'' by ``locally equidimensional.''
\oldpage[IV]{357}
On line 5 of
p.~191, add:
\begin{quote}
If, in addition, the fibres $f^{-1}(f(x))$ of $f$ are equidimensional
($\mathbf{0}_{\mathrm{IV}}$, 14.1.3), then $f$ is said to be
\emph{equidimensional}.
\end{quote}
On the sixth line from the bottom of p.~194, replace
``equidimensional'' by ``locally equidimensional.''

\paragraph{(Err\textsubscript{\textmd{IV}}, 35).}
In the title of (\textbf{IV}, 13.3), on the fourth line from the bottom of p.~194,
replace ``equidimensional'' by ``locally equidimensional.''  In
(\textbf{IV}, 13.3.2), on the sixth line from the bottom of p.~196, make the same
replacement twice.  On the fifth line from the bottom of p.~196, add:
\begin{quote}
If, in addition, the fibres $f^{-1}(f(x))$ of $f$ are equidimensional
($\mathbf{0}_{\mathrm{IV}}$, 14.1.3), then $f$ is said to be
\emph{equidimensional}.
\end{quote}
In (\textbf{IV}, 13.3.7), on line 14 of p.~198, replace ``equidimensional'' by
``locally equidimensional.''  In ({\hyperref[IV.13.3.9]{\textbf{IV}, 13.3.9}}), on the third line from
the bottom of p.~198, replace ``equidimensionality'' by ``local
equidimensionality.''  On lines 10, 12, 20, and 24 of p.~199, replace
``equidimensional'' by ``locally equidimensional.''

\paragraph{(Err\textsubscript{\textmd{IV}}, 36).}
In the introduction to \S14, on the sixth line from the bottom of
p.~199, replace ``equidimensional'' by ``locally equidimensional.''

\paragraph{(Err\textsubscript{\textmd{IV}}, 37).}
In (\textbf{IV}, 14.4.1), on the fifth line from the bottom of p.~209, replace
``locally of finite type'' by ``locally of finite presentation.''  The
present statement of ({\hyperref[IV.14.4.1]{\textbf{IV}, 14.4.1}}) is false even when $f:X\to Y$ is a
closed immersion, $X$ consists of a single closed point $y$ of $Y$, and
$\sh{O}_{Y,y}$ is a field.  Indeed, there exist rings $A$ such that
$Y=\operatorname{Spec}(A)$ is compact, has no isolated points, and is
totally disconnected, while every prime ideal $\mathfrak p$ of $A$ is
maximal and $A_\mathfrak p$ is a field (Bourbaki,
\emph{Alg. comm.}, chap.~II, \S4, exercises 16 and 17).  Clearly the
closed immersion $\{y\}\to Y$ is not open, although it satisfies
conditions b) and b') of ({\hyperref[IV.14.4.1]{\textbf{IV}, 14.4.1}}).

\paragraph{(Err\textsubscript{\textmd{IV}}, 38).}
In (\textbf{IV}, 14.4.1.3), on line 14 of p.~211, replace ``locally quasi-finite
and dominant'' by ``locally quasi-finite, locally of finite presentation,
and dominant.''

\paragraph{(Err\textsubscript{\textmd{IV}}, 39).}
In (\textbf{IV}, 14.4.2), on the tenth line from the bottom of p.~211, replace
``locally of finite type'' by ``locally of finite presentation.''

\paragraph{(Err\textsubscript{\textmd{IV}}, 40).}
In (\textbf{IV}, 14.4.3), on line 18 of p.~212, replace ``locally of finite type''
by ``locally of finite presentation.''

\paragraph{(Err\textsubscript{\textmd{IV}}, 41).}
In ({\hyperref[IV.14.4.4]{\textbf{IV}, 14.4.4}}), on the fifteenth and sixteenth lines from the bottom of
p.~212, replace ``locally of finite type'' by ``locally of finite
presentation.''  On the tenth and eleventh lines from the bottom,
replace ``equidimensional'' by ``locally equidimensional.''

\paragraph{(Err\textsubscript{\textmd{IV}}, 42).}
In (\textbf{IV}, 14.4.8), after line 1 of p.~214, add:
\begin{quote}
If, in addition, the set of irreducible components of $Y$ is locally
finite, these two conditions are equivalent to the following ones:
\end{quote}

\paragraph{(Err\textsubscript{\textmd{IV}}, 43).}
In (\textbf{IV}, 14.4.10), on line 15 of p.~215, replace ``arbitrary'' by ``whose
set of irreducible components is locally finite.''  On line 16 of
p.~215, replace ``of finite type'' by ``of finite presentation.''

\paragraph{(Err\textsubscript{\textmd{IV}}, 44).}
In (\textbf{IV}, 15.5.6.1), on the fifteenth line from the bottom of p.~234,
after the first sentence, add:

\begin{quote}
We can restrict ourselves to the case where in $X$ every irreducible closed subset
has only \emph{one} generic point. Indeed, in the contrary case we
consider the space $Y=X/R$,
\oldpage[IV]{358}
the quotient of $X$ by the equivalence relation
$\overline{\{x\}}=\overline{\{x'\}}$. It follows at once from this
definition that every open set of $X$ is \emph{saturated} for $R$, hence the topology
of $X$ is the \emph{inverse image} of that of $Y$ under the canonical
map $\pi:X\to X/R$. A dense open set in $X$ is of the form
$\pi^{-1}(U)$, where $U$ is a dense open set in $Y$, hence a nowhere dense closed set
$Z$ in $X$ is of the form $\pi^{-1}(S)$, where $S=\pi(Z)$ is a nowhere dense
closed set in $Y$, and it amounts to the same to say that $X\setmin Z$ is connected or
that $Y\setmin S$ is connected. Likewise $\pi(T_x)$ is the set of the
generizations of $\pi(x)$, and it is equivalent to say that
$T_x\setmin\{x\}$ is connected or that $\pi(T_x)\setmin\{\pi(x)\}$ is. Finally, by
construction $Y$ is a locally Noetherian Kolmogorov space,
every irreducible closed subset of which therefore has only one generic point. We see
then that it amounts to the same to prove the lemma for $X$ or for $Y$,
and we can therefore suppose that every irreducible closed subset of $X$ has
a unique generic point.
\end{quote}

\paragraph{(Err\textsubscript{\textmd{IV}}, 45).}
In (\textbf{IV}, 15.6.1), on the eighteenth line from the bottom of p.~236,
delete ``locally Noetherian'' (this hypothesis is not used in the
proof).

\paragraph{(Err\textsubscript{\textmd{IV}}, 46).}
In (\textbf{IV}, 15.6.3), on line 11 of p.~237, replace ``suppose that for every
$y\in Y$'' by ``suppose that $f$ is locally of finite presentation and
that for every $y\in Y$.''  Replace the proof (lines 14--24 of p.~237)
by the following:

\begin{quote}
By {\hyperref[subsection:IV.1.10]{(1.10.3)}}, it is enough to prove that for every $x\in X_y^0$ and
every generalization $y'$ of $y$, there exists a generalization $x'$ of
$x$ such that $f(x')=y'$.  By {\hyperref[IV.15.6.2]{(15.6.2)}}, it is enough to take for $x'$
the generic point of $X_{y'}^0$.
\end{quote}

\paragraph{(Err\textsubscript{\textmd{IV}}, 47).}
In (\textbf{IV}, 15.6.7), on line 2 of p.~241, after ``Noetherian,'' add ``or if
condition $\beta$) of ({\hyperref[IV.15.6.6]{15.6.6}}, (iii)) is satisfied.''  After line 9 of
p.~241, add:
\begin{quote}
When condition $\beta$) of ({\hyperref[IV.15.6.6]{15.6.6}}, (iii)) is satisfied, the implication
a)$\Rightarrow$b') follows from ({\hyperref[IV.15.6.6]{15.6.6}}, (ii)), and the implication
b')$\Rightarrow$a) follows from ({\hyperref[IV.15.6.6]{15.6.6}}, (iii)).
\end{quote}

\paragraph{(Err\textsubscript{\textmd{IV}}, 48).}
In (\textbf{IV}, 20.1.5), line 13 of p.~228; (\textbf{IV}, 20.1.6), line 17 of p.~228;
and (\textbf{IV}, 20.1.7), the thirteenth line from the bottom of p.~228, delete
``strictly.''

\paragraph{(Err\textsubscript{\textmd{IV}}, 49).}
In (\textbf{IV}, 20.2.13), after line 6 of p.~236, add:
\begin{quote}
In particular, if $X$ is reduced and has only finitely many irreducible
components, the notion of torsion-free $\sh{O}_X$-Module defined in
({\hyperref[IV.20]{\textbf{IV}, 20.1.5}}) coincides with that of (\textbf{I}, 7.4.1).
\end{quote}

\paragraph{(Err\textsubscript{\textmd{IV}}, 50).}
In (\textbf{IV}, 20.2.14), on lines 18 and 25 of p.~236, delete ``strictly.''

\paragraph{(Err\textsubscript{\textmd{IV}}, 51).}
In (\textbf{IV}, 20.6.2), on the thirteenth line from the bottom of p.~252,
delete ``strictly.''

\paragraph{(Err\textsubscript{\textmd{IV}}, 52).}
In (\textbf{IV}, 20.6.4), on lines 23 and 22 of p.~253, delete ``strictly.''

\paragraph{(Err\textsubscript{\textmd{IV}}, 53).}
In (\textbf{IV}, 21.4), after line 23 of p.~267, add:

\def\EGAIVSuppressErratumLabels{1}

The following proposition and its corollaries were pointed out to us by
M.~Raynaud.

\begin{proposition}[21.4.9]
\ifdefined\EGAIVSuppressErratumLabels\else\label{IV.21.4.9}\fi
Let $X$, $Y$ be two locally Noetherian preschemes,
$f:X\to Y$ a faithfully flat morphism; suppose moreover \textup{either} that
$f$ is quasi-compact, \textup{or} that $f$ is open. For a divisor
$D'$ on $X$ to be of the form $f^*(D)$, where $D$ is a divisor on
$Y$, it is necessary and sufficient that for every $y\in Y$ such that
$\operatorname{prof}(\sh{O}_{Y,y})\leq 1$, on putting
$Y_1=\operatorname{Spec}(\sh{O}_{Y,y})$,
$X_1=X\times_Y Y_1$,
$f_1=f_{(Y_1)}:X_1\to Y_1$,
the inverse image $D_1'$ of $D'$ on $X_1$ is of the form
$f_1^*(D_1)$, where $D_1\in\operatorname{Div}(Y_1)$.
\end{proposition}

\oldpage[IV]{359}
\begin{proof}
The necessity of the condition is immediate from the transitivity of
inverse images of divisors, all the morphisms in question being flat
{\hyperref[IV.21.4.4]{(21.4.4)}}.  For every open subset $U$ of $Y$, write $f_U$ for the
restriction $f^{-1}(U)\to U$ of $f$.  By {\hyperref[IV.21.4.8]{(21.4.8)}}, the map
\[
 D_U\leadsto f_U^*(D_U)
\]
from $\operatorname{Div}(U)$ to
$\operatorname{Div}(f^{-1}(U))$ is injective.  Thus, if $U_1,U_2$ are
two open subsets of $Y$ and
\[
 D'|f^{-1}(U_1)=f_{U_1}^*(D_1),\qquad
 D'|f^{-1}(U_2)=f_{U_2}^*(D_2),
\]
then necessarily
$D_1|U_1\cap U_2=D_2|U_1\cap U_2$.  It follows at once that there is a
\emph{largest} open subset $U\subset Y$ for which
\[
 D'|f^{-1}(U)=f_U^*(D_U)
 \quad\text{with}\quad D_U\in\operatorname{Div}(U).
\]
We must show that $U=Y$.  Suppose otherwise, and let $y$ be a maximal
point of $Y\setmin U$.  We shall find an open neighbourhood $V$ of
$y$ in $Y$ such that
$D'|f^{-1}(V)=f_V^*(D_V)$ for some
$D_V\in\operatorname{Div}(V)$, contradicting the maximality of $U$.

\textup{I)} \emph{The case where $f$ is quasi-compact.}
Applying the fact that $f$ is quasi-compact and {\hyperref[IV.20.3.8]{(20.3.8)}} following the method of
({\hyperref[IV.8.1.2]{8.1.2}}, \emph{a)}), we are reduced to proving---in the notation of the
statement, but without initially imposing a bound on
$\operatorname{prof}(\sh{O}_{Y,y})$---that $D_1'$ is of the form
$f_1^*(D_1)$.  If $\operatorname{prof}(\sh{O}_{Y,y})\leq 1$, this is
precisely the hypothesis.  Suppose therefore that
$\operatorname{prof}(\sh{O}_{Y,y})\geq 2$.  Since $y$ is a maximal
point of $Y\setmin U$, the open subset $W=Y_1\cap U$ of $Y_1$ is
$Y_1\setmin\{y\}$.  We may therefore restrict to the case
$Y=Y_1$ and $U=Y\setmin\{y\}$.

By hypothesis there is a divisor $D_U\in\operatorname{Div}(U)$ such
that $f_U^*(D_U)=D'|f^{-1}(U)$.  The divisor $D_U$ is the equivalence
class of a pair $(\sh{L}_U,s)$, where $\sh{L}_U$ is an invertible
$\sh{O}_U$-Module and $s$ a regular meromorphic section of $\sh{L}_U$
over $U$ {\hyperref[IV.21.2.11]{(21.2.11)}}.  Hence, by {\hyperref[IV.21.4.1]{(21.4.1)}}, $D'|f^{-1}(U)$ is the
equivalence class of $(\sh{L}'_U,s')$, where
\[
 \sh{L}'_U=f_U^*(\sh{L}_U),\qquad s'=s\circ f_U.
\]
Let $i:U\to Y$ and $j:f^{-1}(U)\to X$ be the canonical injections,
and put
\[
 \sh{L}=i_*(\sh{L}_U),\qquad
 \sh{L}'=j_*(\sh{L}'_U).
\]
Since $f$ is flat, $f^*(\sh{L})$ is canonically isomorphic to
$\sh{L}'$ {\hyperref[IV.2.3.1]{(2.3.1)}}.  Moreover, flatness and the inequality
$\operatorname{prof}(\sh{O}_{Y,y})\geq 2$ imply
$\operatorname{prof}(\sh{O}_{X,x})\geq 2$ for \emph{every}
$x\in f^{-1}(y)$ {\hyperref[IV.6.3.1]{(6.3.1)}}.

The divisor $D'$ is the equivalence class of a pair
$(\sh{L}'',s'')$, where $\sh{L}''$ is an invertible $\sh{O}_X$-Module
and $s''$ a regular meromorphic section of $\sh{L}''$ over $X$.
Consequently $\sh{L}'_U$ is isomorphic to
$\sh{L}''|f^{-1}(U)$, and the preceding observations together with
{\hyperref[IV.5.9.8]{(5.9.8)}} and {\hyperref[IV.5.10.5]{(5.10.5)}} show that $\sh{L}'$ is an invertible
$\sh{O}_X$-Module isomorphic to $\sh{L}''$.  Since $f$ is faithfully
flat and quasi-compact, it follows that $\sh{L}$ is an invertible
$\sh{O}_Y$-Module {\hyperref[IV.2.5.2]{(2.5.2)}}, and hence isomorphic to $\sh{O}_Y$, because
$Y$ is a local scheme.

Furthermore, since $\operatorname{prof}(\sh{O}_{Y,y})\geq 1$, the open
subset $U=Y\setmin\{y\}$ is schematically dense in $Y$
({\hyperref[IV.20.2.13]{20.2.13}}, (iv)).  The regular meromorphic section $s$ of $\sh{L}_U$
over $U$ therefore extends to a regular meromorphic section $s_0$ of
$\sh{L}$ over $Y$ {\hyperref[IV.20.2.11]{(20.2.11)}}.  The equivalence class of
$(\sh{L},s_0)$ is thus a divisor $D$ on $Y$.  Plainly $f^*(D)$ and
$D'$ induce the same divisor on
$f^{-1}(U)=X\setmin f^{-1}(y)$; and since
$\operatorname{prof}(\sh{O}_{X,x})\geq 2$ for every
$x\in f^{-1}(y)$, {\hyperref[IV.21.1.8]{(21.1.8)}} gives $D'=f^*(D)$.  This proves case I.

\textup{II)} \emph{The case where $f$ is open.}
The problem is to find an open neighbourhood $V$ of $y$ in $Y$ such
that $D'|f^{-1}(V)=f_V^*(D_V)$ for some
$D_V\in\operatorname{Div}(V)$, so we may first suppose $Y$ affine.
Because $f$ is open, the images $f(W)$, with $W$ running through the
affine open subsets of $X$, form an open covering of $Y$.  Since $Y$
is quasi-compact, there is a \emph{quasi-compact} open subset $T\subset X$
such that $f(T)=Y$, and $f|T$ is quasi-compact (1.1.1).  A fortiori,
for every quasi-compact open $T_\alpha\supset T$ in $X$,
\oldpage[IV]{360}
one has
$f(T_\alpha)=Y$ and $f|T_\alpha$ is quasi-compact.

Applying case I to $f|T_\alpha$, we obtain a divisor $D_\alpha$ on
$Y$ such that
\[
 (f|T_\alpha)^*(D_\alpha)=D'|T_\alpha.
\]
By transitivity of inverse images of divisors under flat morphisms
{\hyperref[IV.21.4.4]{(21.4.4)}}, one also has
$(f|T)^*(D_\alpha)=D'|T$.  Hence for any $\alpha,\beta$,
\[
 (f|T)^*(D_\alpha)=(f|T)^*(D_\beta),
\]
and {\hyperref[IV.21.4.8]{(21.4.8)}} gives $D_\alpha=D_\beta$.  If $D$ denotes their common
value, then $f^*(D)|T_\alpha=D'|T_\alpha$ for every $\alpha$, and
therefore $f^*(D)=D'$. C.Q.F.D.
\end{proof}

\begin{corollary}[21.4.10]
\ifdefined\EGAIVSuppressErratumLabels\else\label{IV.21.4.10}\fi
Assume the general hypotheses of {\hyperref[IV.21.4.9]{(21.4.9)}}, and suppose in addition that
$Y$ is normal and that, for every $y\in Y$ with
$\dim(\sh{O}_{Y,y})=1$, the fibre $X_y=f^{-1}(y)$ has no embedded
associated prime cycle.  A divisor $D'$ on $X$ is of the form $f^*(D)$,
with $D\in\operatorname{Div}(Y)$, if and only if the following two
conditions hold:
\begin{enumerate}[label=\textup{\arabic*\textdegree}]
\item for every maximal point $y$ of $Y$, the inverse image of $D'$ in
$\operatorname{Div}(X_y)$ is zero;
\item for every $y\in Y$ such that $\dim(\sh{O}_{Y,y})=1$, there is an
integer $n_y$ such that, for every maximal point $x$ of $X_y$, the
inverse image of $D'$ in
$\operatorname{Div}(\operatorname{Spec}(\sh{O}_{X,x}))$ equals
$\operatorname{div}(t_y^{n_y})$, where $t_y$ is the image under
$\sh{O}_{Y,y}\to\sh{O}_{X,x}$ of a uniformizer of the discrete
valuation ring $\sh{O}_{Y,y}$.
\end{enumerate}
\end{corollary}

\begin{proof}
Apply the criterion {\hyperref[IV.21.4.9]{(21.4.9)}}.  At a maximal point $y$ of $Y$,
$\operatorname{Div}(Y_1)=0$ because $\sh{O}_{Y,y}$ is Artinian;
condition 1\textdegree{} is therefore equivalent to the criterion of
{\hyperref[IV.21.4.9]{(21.4.9)}} at such a point.  It remains to show that conditions
1\textdegree{} and 2\textdegree{} likewise imply the criterion at a
point $y\in Y$ with $\operatorname{prof}(\sh{O}_{Y,y})=1$.  Since $Y$
is normal, these are precisely the points for which
$\dim(\sh{O}_{Y,y})=1$, equivalently those for which
$\sh{O}_{Y,y}$ is a discrete valuation ring {\hyperref[IV.5.8.6]{(5.8.6)}}.

We may thus restrict to $Y=\operatorname{Spec}(\sh{O}_{Y,y})$.  If $t$
is a uniformizer of $\sh{O}_{Y,y}$, the map
$n\leadsto\operatorname{div}(t^n)$ is an isomorphism from $\mathbb Z$
onto $\operatorname{Div}(Y)$.  The condition of {\hyperref[IV.21.4.9]{(21.4.9)}} says that
$D'$ is of the form $\operatorname{div}(t^n\circ f)$ for some
$n\in\mathbb Z$, and plainly implies condition 2\textdegree{}.
Conversely, suppose 2\textdegree{} holds, as does 1\textdegree{}, and
write $n=n_y$.  It remains to prove
$D'=\operatorname{div}(t^n\circ f)$.  If $\eta$ is the generic point of
$Y$, the germs of these two divisors at every point of $X_\eta$ are
zero by condition 1\textdegree{}.  Their germs at every maximal point
of $X_y$ are equal by condition 2\textdegree{}.  At a nonmaximal point
$x\in X_y$, one has
$\operatorname{prof}(\sh{O}_{X_y,x})\geq 1$, because $X_y$ has no
embedded associated prime cycle, and hence
$\operatorname{prof}(\sh{O}_{X,x})\geq 2$ by {\hyperref[IV.6.3.1]{(6.3.1)}}.  The desired
equality follows from {\hyperref[IV.21.1.8]{(21.1.8)}}.
\end{proof}

\begin{corollary}[21.4.11]
\ifdefined\EGAIVSuppressErratumLabels\else\label{IV.21.4.11}\fi
Assume the general hypotheses of {\hyperref[IV.21.4.9]{(21.4.9)}}, and suppose in addition that
$Y$ is normal and that, for every $y\in Y$ with
$\dim(\sh{O}_{Y,y})=1$, the fibre $X_y$ is integral.  A divisor $D'$ is
of the form $f^*(D)$, with $D\in\operatorname{Div}(Y)$, if and only if,
for every maximal point $y$ of $Y$, the inverse image of $D'$ in
$\operatorname{Div}(X_y)$ is zero.
\end{corollary}

\begin{proof}
It is enough to check that condition 2\textdegree{} of {\hyperref[IV.21.4.10]{(21.4.10)}} is then
automatic.  If $\dim(\sh{O}_{Y,y})=1$ and $X_y$ is integral, then at
the unique maximal point $x$ of $X_y$ one has
$\dim(\sh{O}_{X,x})=1$ by {\hyperref[IV.6.1.3]{(6.1.3)}}.  Moreover,
$\sh{O}_{Y,y}$ is regular and $\sh{O}_{X_y,x}$ is a field, so
$\sh{O}_{X,x}$ is regular {\hyperref[IV.6.5.1]{(6.5.1)}}, hence a discrete valuation ring.
If $t$ is a uniformizer of $\sh{O}_{Y,y}$, its image $t_y$ in
$\sh{O}_{X,x}$ is a uniformizer of $\sh{O}_{X,x}$
($\mathbf{0}$, 17.3.3).  The inverse image of $D'$ in
$\operatorname{Div}(\operatorname{Spec}(\sh{O}_{X,x}))$ is therefore
\oldpage[IV]{361}
of the form $\operatorname{div}(t_y^n)$, which is exactly condition
2\textdegree{} of {\hyperref[IV.21.4.10]{(21.4.10)}}, since $X_y$ has only one maximal point.
\end{proof}

\begin{remark}[21.4.12]
\ifdefined\EGAIVSuppressErratumLabels\else\label{IV.21.4.12}\fi
If in {\hyperref[IV.21.4.11]{(21.4.11)}} one assumes only that the $X_y$ are \emph{reduced}, one must
add to the condition in the statement that the inverse images of $D'$
in \emph{each} group
$\operatorname{Div}(\operatorname{Spec}(\sh{O}_{X,x}))$ with $x$
running over the maximal points of $X_y$, are all of the form
$\operatorname{div}(t'^n_x)$, where $t_x'$ is a uniformizer of
$\sh{O}_{X,x}$ and the integer $n$ is \emph{the same} for all these points.
\end{remark}

\begin{corollary}[21.4.13]
\ifdefined\EGAIVSuppressErratumLabels\else\label{IV.21.4.13}\fi
Assume the hypotheses of \textup{\hyperref[IV.21.4.11]{(21.4.11)}}.  For every invertible
$\sh{O}_X$-Module $\sh{L}$ such that, for every maximal point $\eta$ of
$Y$, the inverse image
\[
 \sh{L}_\eta
 =\sh{L}\otimes_{\sh{O}_X}\sh{O}_{X_\eta}
\]
is isomorphic to $\sh{O}_{X_\eta}$, there exist an invertible
$\sh{O}_Y$-Module $\sh{M}$ and an isomorphism
$f^*(\sh{M})\simeq\sh{L}$.
\end{corollary}

\begin{proof}
It is enough to find a regular meromorphic section $s$ of $\sh{L}$ over
$X$ such that the divisor $D'$ corresponding to $(\sh{L},s)$ has
$D'_x=0$ at every point $x$ of a fibre $X_\eta$ over a maximal point
$\eta$ of $Y$.  Applying {\hyperref[IV.21.4.11]{(21.4.11)}} to $D'$ then gives the result by
{\hyperref[IV.21.4.2]{(21.4.2)}}.

The points of $\operatorname{Ass}(\sh{O}_X)$ lie over the maximal
points of $Y$ {\hyperref[IV.3.3.1]{(3.3.1)}}.  Since $Y$ is reduced and locally Noetherian,
there is a dense open subset $V$ of $Y$ which is the union of a
disjoint family of affine open subsets, each containing exactly one
maximal point of $Y$; and $f^{-1}(V)$ is schematically dense in $X$
{\hyperref[IV.5.10.2]{(5.10.2)}}.  We may therefore restrict to the case in which $Y$ is
affine and integral.

For the unique maximal point $\eta$ of $Y$, the fibre $X_\eta$ is then
integral, so $\operatorname{Ass}(\sh{O}_X)$ consists of a single point
$z$, the generic point of $X$.  Let $s_\eta$ be the section of
$\sh{L}_\eta$ over $X_\eta$ which corresponds, under an isomorphism,
to the unit section of $\sh{O}_{X_\eta}$.  Let $U$ be an affine open
neighbourhood in $X$ of a point $x\in X_\eta$, and let $s$ be a section
of $\sh{L}$ over $U$ with the same germ at $x$ as $s_\eta$.  Since $U$
is schematically dense in $X$ ({\hyperref[IV.20.2.13]{20.2.13}}, (iv)), $s$ identifies with a
uniquely determined regular meromorphic section of $\sh{L}$ over $X$.
It is clear from {\hyperref[IV.20.2.11]{(20.2.11)}} and {\hyperref[IV.20]{(20.3.5)}} that this section induces on
$X_\eta$ a regular meromorphic section of $\sh{L}_\eta$ equal to
$s_\eta$, and hence has the required property.
\end{proof}

\let\EGAIVSuppressErratumLabels\undefined
}

\oldpage[IV]{349}
\begin{env}[3.1.10.5]
\label{IV.3.1.10.5}
We note that even when $f$ is a proper morphism, one does not necessarily have
$f(\operatorname{Ass}(\sh{F}))\subset\operatorname{Ass}(f_*(\sh{F}))$:
one sees this by taking for $Y$ the spectrum of a field and for $X$ a
projective space over $k$, for example $\mathbf P^1_k$; then
$\sh{F}=\sh{O}_X(-1)$ is a coherent $\sh{O}_X$-Module $\neq0$ such
that $\Gamma(X,\sh{F})=0$ (\textbf{III}, 2.1.13), hence
$\operatorname{Ass}(f_*(\sh{F}))=\varnothing$ and
$\operatorname{Ass}(\sh{F})\neq\varnothing$.
\end{env}

\paragraph{($\mathbf{Err}_{\mathrm{IV}}$, 18)}
After (\textbf{IV}, 3.2.8), insert after line 17 of p.~43:

\begin{remark}[3.2.9]
\label{IV.3.2.9}
For every point $s$ of a prescheme $X$, write
$i_s:\operatorname{Spec}(\boldsymbol{k}(s))\to X$ for the canonical morphism, which
is a quasi-compact monomorphism.  For every $\boldsymbol{k}(s)$-module $G(s)$,
regarded as a Module on $\operatorname{Spec}(\boldsymbol{k}(s))$,
$(i_s)_*(G(s))$ is therefore a quasi-coherent $\sh{O}_X$-Module
{\hyperref[subsection:IV.1.7]{(1.7.4)}}.  Its fibre is zero at every point
$x\notin\overline{\{s\}}$ and is isomorphic to $G(s)$ otherwise.  If
$X$ is locally Noetherian and $G(s)$ has finite length, then
$\operatorname{Ass}((i_s)_*(G(s)))$ consists of the single point $s$;
in other words, $(i_s)_*(G(s))$ is irredundant.

Now suppose that $X$ is locally Noetherian and that $\sh{F}$ is a
coherent $\sh{O}_X$-Module.  Let $S$ be a discrete subset of $X$, and
for every $s\in S$ put $\sh{F}(s)=i_s^*(\sh{F})$ a $\boldsymbol{k}(s)$-module.
For each $s\in S$, choose a finite-length quotient $G(s)$ of
$\sh{F}(s)$.  There is a canonical homomorphism
\[
 u:\sh{F}\longrightarrow
 \bigoplus_{s\in S}(i_s)_*(\sh{F}(s))\longrightarrow
 \bigoplus_{s\in S}(i_s)_*(G(s))
\]
of coherent $\sh{O}_X$-Modules, since a discrete subset of $X$ is
locally finite.  If $\sh{H}=\operatorname{Ker}u$ and
$\sh{G}=\sh{F}/\sh{H}$, then $\sh{G}$ is coherent and $u$ induces an
injection
\[
 v:\sh{G}\longrightarrow\bigoplus_{s\in S}(i_s)_*(G(s)).
\]
If $\sh{G}^{(s)}$ denotes the image of $\sh{G}$ in
$(i_s)_*(G(s))$, the family $(\sh{G}^{(s)})_{s\in S}$ is an
irredundant decomposition, reduced when all the $G(s)$ are nonzero, and
$\operatorname{Ass}(\sh{G})=S$; thus $\sh{G}$ has no embedded
associated prime cycle.

This gives a bijection $(G(s))_{s\in S}\leadsto\sh{G}$, where
$S$ ranges over the discrete subsets of $X$ for which
$\sh{F}(s)\neq0$ for every $s\in S$, and, for each $S$, the $G(s)$
range over the nonzero finite-length quotients of $\sh{F}(s)$; on the
other side, $\sh{G}$ ranges over the coherent quotient
$\sh{O}_X$-Modules of $\sh{F}$ without embedded associated prime
cycles.  The inverse bijection associates to $\sh{G}$ its reduced
irredundant decomposition.
\end{remark}

\paragraph{($\mathbf{Err}_{\mathrm{IV}}$, 19)}
In (\textbf{IV}, 4.2.1), before the seventeenth line from the bottom of p.~55,
insert:

\begin{remark}[4.2.1.4]
\label{IV.4.2.1.4}
Under the hypotheses of {\hyperref[IV.4.2.1]{(4.2.1)}}, we have
moreover
\begin{equation}
 \dim(K\otimes_kL)=
 \inf(\operatorname{tr.deg}_kK,\operatorname{tr.deg}_kL)
 \label{IV.4.2.1.5}\tag{4.2.1.5}
\end{equation}
where it is to be understood that if the two fields $K$, $L$ are of infinite
transcendence degree (which may be arbitrary \emph{cardinals}), the
left-hand side is to be taken equal to $+\infty$ (the dimension considered
here is the one defined in (\textbf{0}, 16.1.1)).

Suppose indeed, for example, that
$\operatorname{tr.deg}_kK\leq\operatorname{tr.deg}_kL$, and in the first
place that $\operatorname{tr.deg}_kK=n<+\infty$. We have seen in the proof
of {\hyperref[IV.4.2.1]{(4.2.1)}} that $K\otimes_kL$ is an
algebraic extension of $K'\otimes_kL$, and therefore has the same dimension
(\textbf{0}, 16.1.5). We can therefore restrict ourselves to the case where
$K=k(T_1,\ldots,T_n)$; then $K\otimes_kL$ is a ring of fractions
of
\oldpage[IV]{350}
$k[T_1,\ldots,T_n]\otimes_kL=L[T_1,\ldots,T_n]$,
which is of dimension $n$ {\hyperref[IV.5.5.4]{(5.5.4)}}, and
consequently $\dim(K\otimes_kL)\leq n$ (\textbf{0}, 16.1.3.1). To prove that
$\dim(K\otimes_kL)=n$, it obviously suffices to show that there is in
$\operatorname{Spec}(L[T_1,\ldots,T_n])$ a point rational over $L$
which belongs to $\operatorname{Spec}(K\otimes_kL)$
{\hyperref[IV.5.2.3]{(5.2.3)}}, or again (\textbf{I}, 1.7.4) that there
exists an $L$-\emph{homomorphism} $u:K\otimes_kL\to L$. But by definition
of the tensor product, this amounts to proving the existence of a
$k$-\emph{homomorphism} $v:K\to L$, which is evident by virtue of
the hypothesis
$\operatorname{tr.deg}_kK\leq\operatorname{tr.deg}_kL$.
If $\operatorname{tr.deg}_kK\leq\operatorname{tr.deg}_kL$ and if these
two cardinals are infinite, one sees in the same way that there is a
rational point in $\operatorname{Spec}(L[\mathbf{t}])$ which
belongs to $\operatorname{Spec}(K\otimes_kL)$; but the kernel of every
$L$-homomorphism from $L[\mathbf{t}]$ to $L$ is a
maximal ideal generated by a family of the form
$(t_\alpha\setmin a_\alpha)_{\alpha\in I}$ with $a_\alpha\in L$. It is clear
that such a maximal ideal contains an infinite decreasing sequence
of distinct prime ideals, hence the corresponding local ring of
$L[\mathbf{t}]$ is non-Noetherian and of infinite
dimension; since it is a local ring of $K\otimes_kL$, we
see that the latter is non-Noetherian and of infinite dimension.
\end{remark}

\paragraph{($\mathbf{Err}_{\mathrm{IV}}$, 20)}
After (\textbf{IV}, 4.5.19), before the twelfth line from the bottom of p.~67,
insert:

\begin{corollary}[4.5.19.3]
\label{IV.4.5.19.3}
Let $k$ be a field and $X$ a finite-type $k$-prescheme which is
irreducible but not geometrically irreducible.  There is a closed subset
$Y$ of $X$, nowhere dense in $X$, containing every $x\in X$ for which
$\boldsymbol{k}(x)$ is radicial over $k$ (and in particular every $k$-rational point
of $X$).
\end{corollary}

Let $k'$ be an algebraic closure of $k$, let $X'=X\otimes_kk'$, and let
$p:X'\to X$ be the canonical projection.  If $\boldsymbol{k}(x)$ is radicial over
$k$, there is only one point $x'\in X'$ over $x$ (\textbf{I}, 3.4.9); this point
must belong to two distinct irreducible components of $X'$, since
otherwise $X$ would be geometrically irreducible by {\hyperref[IV.4.5.19]{(4.5.19)}}.  Let
$X'_i$ ($1\leq i\leq n$) be the irreducible components of the Noetherian
prescheme $X'$, and for $1\leq i<j\leq n$ put
$X'_{ij}=X'_i\cap X'_j$.  Then
\[
 Y=p\left(\bigcup_{i<j}X'_{ij}\right)
\]
has the required properties.  Since $p$ is integral and the $X'_{ij}$
are closed in $X'$, $Y$ is closed in $X$.  If $\xi$ is the generic point
of $X$, $p^{-1}(\xi)$ is discrete and consists of the generic points of
the $X'_i$ {\hyperref[IV.4.5.11]{(4.5.11)}}, so it meets none of the $X'_{ij}$; hence $Y$ is
nowhere dense.

\paragraph{($\mathbf{Err}_{\mathrm{IV}}$, 21)}
In (\textbf{IV}, 5.6.11.1), on the seventh line from the bottom of p.~102, add
``of dimension $2$'' after ``integral.''

\paragraph{($\mathbf{Err}_{\mathrm{IV}}$, 22)}
After (\textbf{IV}, 6.2.3), after line 18 of p.~138, insert:

\begin{proposition}[6.2.4]
\label{IV.6.2.4}
Under the hypotheses of {\hyperref[IV.6.2.3]{(6.2.3)}}, let in addition $M$ be a finite
$A$-module, and suppose that $M\neq0$ and $N\neq0$.  Then
\begin{equation}
 \operatorname{pd}_B(M\otimes_AN)=
 \operatorname{pd}_A(M)+
 \operatorname{pd}_{B\otimes_Ak}(N\otimes_Ak). \label{IV.6.2.4.1}
\tag{6.2.4.1}
\end{equation}
\end{proposition}

Consider the left resolution $L_\bullet$ of $N$ introduced in {\hyperref[IV.6.2.3]{(6.2.3)}},
and a left resolution $K_\bullet$ of $M$ by finite free $A$-modules:
\[
 \ldots\longrightarrow K_i\longrightarrow K_{i-1}\longrightarrow
 \ldots\longrightarrow K_0\longrightarrow M\longrightarrow0
\]
We saw that the $L_i$, $N$, and the $Z_i(L_\bullet)$ are flat
$A$-modules.  The same is true of the $B_i(L_\bullet)$, which equal
$Z_i(L_\bullet)$ except for $i=0$; in that case $B_0(L_\bullet)$ is flat
\oldpage[IV]{351}
by ($\mathbf{0}_{\mathrm I}$, 6.1.2).  Finally the $H_i(L_\bullet)$ vanish except
for $H_0(L_\bullet)=N$, which is flat over $A$.  Under these conditions
one knows that
$H_n(K_\bullet\otimes_AL_\bullet)$ is canonically isomorphic to
$\bigoplus_{p+q=n}(H_p(K_\bullet)\otimes_AH_q(L_\bullet))$ (G, I, 5.5).  Thus $K_\bullet\otimes_AL_\bullet$, with the usual total
grading, is a resolution of $M\otimes_AN$ by free $B$-modules.

Put
\[
 r=\operatorname{pd}_AM,\qquad
 s=\operatorname{pd}_B(N)=
   \operatorname{pd}_{B\otimes_Ak}(N\otimes_Ak).
\]
There are free resolutions $K_\bullet$ and $L_\bullet$ of lengths $r$
and $s$, so their tensor product is a free resolution of length $r+s$;
this proves
$\operatorname{pd}_B(M\otimes_AN)\leq r+s$.  It remains, by
(\textbf{0}, 17.2.6), to show
\[
 \operatorname{Tor}^{B}_{r+s}(M\otimes_AN,\boldsymbol{k}(B))\neq0.
\]
Using the preceding resolution, the left-hand side is
\[
 H_{r+s}((K_\bullet\otimes_AL_\bullet)\otimes_B\boldsymbol{k}(B)).
\]
But
\[
 (K_\bullet\otimes_AL_\bullet)\otimes_B\boldsymbol{k}(B)
 = (K_\bullet\otimes_Ak)\otimes_k
   (L_\bullet\otimes_B\boldsymbol{k}(B)).
\]
This is a tensor product of complexes of vector spaces over $k$, so the
K\"unneth formula gives
\[
 H_{r+s}((K_\bullet\otimes_Ak)\otimes_k(L_\bullet\otimes_B\boldsymbol{k}(B)))=H_r(K_\bullet\otimes_Ak)\otimes_k
 H_s(L_\bullet\otimes_B\boldsymbol{k}(B)),
\]
since the homology of the first, respectively second, complex vanishes
above $r$, respectively $s$.  Hence, up to canonical isomorphism,
\[
 \operatorname{Tor}^{B}_{r+s}(M\otimes_AN,\boldsymbol{k}(B))=
 \operatorname{Tor}^{A}_{r}(M,k)\otimes_k
 \operatorname{Tor}^{B}_{s}(N,\boldsymbol{k}(B)),
\]
and both factors are nonzero by hypothesis; so is their tensor product.

\oldpage[IV]{352}
\paragraph{(Err\textsubscript{\textmd{IV}}, 23).}
In (\textbf{IV}, 6.10.1), on line 8 of p.~155, add after ``canonical'':
``$\sh{F}$ a coherent $\sh{O}_X$-Module.''

\paragraph{(Err\textsubscript{\textmd{IV}}, 24).}
In ({\hyperref[IV.6.14.4]{\textbf{IV}, 6.14.4}}), the reference on the twelfth line from the bottom of
p.~175 is insufficient; argue as follows.  First, if $M$ is a
transcendence basis of $K'$ over $A$, then $A(M)$ is algebraically closed
in $B(A(M))$ (Bourbaki, \emph{Alg.}, chap.~V, \S6, exercise 8(c)).  We
may therefore reduce to the case in which $K'$ is an algebraic separable
extension of $A$.

Let $p$ be the characteristic exponent of $A$.  It follows from
Bourbaki, \emph{loc. cit.}, \S9, exercise 2, that the algebraic closure
of $K'$ in $B(K')$ is radicial over $K'$.  We show that every
$x\in B(K')$ such that $x^{p^r}\in K'$ already belongs to $K'$.
Choose a basis $(e_\lambda)$ of $K'$ over $A$.  Since $K'/A$ is
algebraic separable, the $e_\lambda^{p^r}$ also form a basis of $K'$ over
$A$ (Bourbaki, \emph{loc. cit.}, \S8, no.~3, corollary to prop.~4).
Write $x=\sum_\lambda e_\lambda b_\lambda$ with $b_\lambda\in B$.  By
hypothesis there are $a_\lambda\in A$ such that
\[
 \sum_\lambda e_\lambda^{p^r}b_\lambda^{p^r}
 =\sum_\lambda e_\lambda^{p^r}a_\lambda.
\]
Therefore $b_\lambda^{p^r}=a_\lambda$ for every $\lambda$.  Since $A$
is algebraically closed in $B$, each $b_\lambda$ lies in $A$, and hence
$x\in K'$.

\paragraph{(Err\textsubscript{\textmd{IV}}, 25).}
In (\textbf{IV}, 7.4.8, B), p.~203, the question was answered affirmatively by
R.~\textsc{Kiehl}, ``Ausgezeichnete Ringe der nichtarchimedischen analytischen
Geometrie,'' forthcoming in \emph{Inventiones Mathematicae}.

\paragraph{(Err\textsubscript{\textmd{IV}}, 26).}
In (\textbf{IV}, 7.8.4), on the ninth line from the bottom of p.~217, replace
``We shall see later (\S18) that if $k$ is a complete nondiscrete valued
field\ldots'' by:
\begin{quote}
R.~\textsc{Kiehl} proved (in the article cited in
\textbf{(Err\textsubscript{\textmd{IV}}, 25)}) that if $k$ is a complete nondiscrete
valued field\ldots
\end{quote}

\paragraph{(Err\textsubscript{\textmd{IV}}, 27).}
In (\textbf{IV}, 7.8.4), insert before the sixth line from the bottom of p.~217:

\smallskip
\noindent
\textup{(vi)} Let $A$ be a one-dimensional Noetherian local integral
ring.  It is then universally catenary {\hyperref[IV.7.2.9]{(7.2.9)}}, and hence, by {\hyperref[IV.7.3.19]{(7.3.19)}},
it is excellent if and only if its formal fibres are geometrically
reduced; equivalently, if and only if its completion $\widehat A$ is
reduced and the component fields of the total ring of fractions of
$\widehat A$ are separable extensions of the fraction field $K$ of
$A$.  By {\hyperref[IV.7.6.4]{(7.6.4)}} and {\hyperref[IV.7.7.2]{(7.7.2)}}, this is also equivalent to saying that
$A$ is universally Japanese.  When $A$ is a discrete valuation ring,
so is $\widehat A$; to say that $A$ is excellent then means that the
fraction field of $\widehat A$ is a separable extension of $K$, a
condition which is always satisfied when $K$ has characteristic zero.

\paragraph{(Err\textsubscript{\textmd{IV}}, 28).}
After (\textbf{IV}, 8.10.3), following line 14 of p.~37, insert:

\begin{remark}[8.10.3.1]
\label{IV.8.10.3.1}
The hypothesis in {\hyperref[IV.8.10.3]{(8.10.3)}} that the sets
$\overline{f_\alpha(X_\alpha)}$ be constructible cannot be omitted.
For example, take $S=\operatorname{Spec}(A)$, where $A$ is a ring such
that $S$ is compact, totally disconnected, and has no isolated point,
every prime ideal $\mathfrak p$ of $A$ is maximal, and $A_\mathfrak p$
is a field (Bourbaki, \emph{Alg. comm.}, chap.~II, \S4,
exercises 16 and 17).  Let $s\in S$, and take for
$Y_\alpha=S_\alpha$ the affine open neighbourhoods of $s$ in $S$.
For every $\alpha$, take $X_\alpha=\operatorname{Spec}(\boldsymbol{k}(s))$, with
$f_\alpha$ and $f$ the canonical morphisms.  It is clear that $f$,
which is the identity, is dominant, whereas none of the $f_\alpha$ is.
\end{remark}

\paragraph{(Err\textsubscript{\textmd{IV}}, 29).}
In (\textbf{IV}, 8.12.8), on line 7 of p.~47, replace ``such that there exists an
ample $\sh{O}_Y$-Module'' by ``and quasi-separated.''  On line 12 of
p.~47, replace ``The hypothesis'' by:
\begin{quote}
It follows from {\hyperref[IV.8.12.3]{(8.12.3)}} that the existence of the factorisation in the
statement is a property local on $Y$; we may therefore suppose $Y$
\emph{affine}, which
\end{quote}

\paragraph{(Err\textsubscript{\textmd{IV}}, 30).}
After (\textbf{IV}, 8.12.11), add at the end of p.~48:

\begin{corollary}[8.12.12]
\label{IV.8.12.12}
Let $X$, $Y$ be two integral preschemes,
$f:X\to Y$ a separated, birational morphism locally of finite type.
Let $y\in Y$ be a normal point of $Y$, and suppose that a point
$x\in f^{-1}(y)$ is isolated in $f^{-1}(y)$. Then there exists an
open neighbourhood $U$ of $y$ in $Y$ such that the restriction
$f^{-1}(U)\to U$ of $f$ is an isomorphism.
\end{corollary}

\begin{proof}
It suffices to prove that there exist an open neighbourhood $U$ of $y$ in
$Y$ and a $U$-section $g:U\to f^{-1}(U)$ of $f^{-1}(U)$ such that
$g(y)=x$. Indeed, since $U$ contains the generic point $\eta$ of
$Y$ and $f^{-1}(U)$ the generic point $\xi$ of $X$, we necessarily
have $g(\eta)=\xi$ since $f$ is birational
{\hyperref[IV.6.15.4]{(6.15.4)}}; since $g$ is an
immersion, $g(U)$ is locally closed in $f^{-1}(U)$, hence open
in its closure in $f^{-1}(U)$, and since the latter is equal
to $f^{-1}(U)$ by what precedes, $g(U)$ is \emph{open} in
$f^{-1}(U)$; since the sub-prescheme of $f^{-1}(U)$ associated with $g$ is
isomorphic to $U$, hence reduced, and $f^{-1}(U)$ is reduced, it follows
from (\textbf{I}, 5.2.1) that $g$ is an \emph{open immersion}. On the other hand, since
$f$ is separated, $g$ is a \emph{closed immersion} (\textbf{I}, 5.4.6), hence $g(U)$
is a set both open and closed in $f^{-1}(U)$. But since
$f^{-1}(U)$ is integral, this will imply $g(U)=f^{-1}(U)$, whence the
conclusion.

Let us then consider two cases:

\textup{I)} Suppose first that $Y$ is \emph{normal}, and let $V$ be the set of
points of $X$ at which $f$ is quasi-finite; it is an \emph{open} set
in $X$ {\hyperref[IV.13.1.4]{(13.1.4)}}\footnote{The
reader will verify that neither {\hyperref[IV.8.12.12]{(8.12.12)}}
nor any of its consequences is used in the proof of
{\hyperref[IV.13.1.4]{(13.1.4)}}.}. It then follows from
{\hyperref[IV.8.12.10]{(8.12.10)}}
\oldpage[IV]{353}
that the restriction $f|V:V\to Y$ is an open immersion; hence there is
an open neighbourhood $U$ of $y\in Y$ such that $f|V$ is an isomorphism
of $V$ onto the sub-prescheme induced by $Y$ on $U$; the inverse
isomorphism $g$ is then a $U$-section of $f^{-1}(U)$ such that
$g(y)=x$.

\textup{II)} Obtaining a $U$-section $g$ such that $g(y)=x$ being
a problem local on $X$ and on $Y$, we may suppose that $X$ and $Y$
are affine and that $f$ is of finite type, hence that $X$ is a
closed sub-prescheme of a $Y$-prescheme $Z$ of finite presentation over
$Y$ (for example of the form $Y[T_1,\ldots,T_r]$). Put $Y'=\operatorname{Spec}(\sh{O}_{Y,y})$, $X'=X_{(Y')}$, $Z'=Z_{(Y')}$,
so that $X'$ is a closed sub-prescheme of $Z'$. Moreover, if
$f'=f_{(Y')}$, we have $f'^{-1}(y)=f^{-1}(y)$ (\textbf{I}, 3.6.5), and $x$ is therefore
isolated in $f'^{-1}(y)$. Since $Y'$ is \emph{normal}, we can apply the
result of I) to $f'$, and since $Y'$ is the only neighbourhood of $y$ in
$Y'$, there exists a $Y'$-\emph{section} $g'$ \emph{of} $X'$ such that $g'(y)=x$, which
can obviously also be considered as a $Y'$-section \emph{of} $Z'$.
Since the structural morphism $h:Z\to Y$ is of finite presentation,
we can apply ({\hyperref[IV.8.8.2]{8.8.2}}, (i) and
(ii)) following the method of
({\hyperref[IV.8.1.2]{8.1.2}}, \emph{a)}), whence we conclude
the existence of an open neighbourhood $U$ of $y$ in $Y$ and of a
$U$-section $g$ \emph{of} $h^{-1}(U)$ such that $g(y)=x$. It remains to see
that in fact $g$ is a $U$-section \emph{of} $f^{-1}(U)$. But since
$f^{-1}(U)$ is a closed sub-prescheme of $h^{-1}(U)$ and $U$ is
reduced, it suffices for this (\textbf{I}, 5.2.2) to see that
$g(U)\subset f^{-1}(U)$. Now, since $U$ is integral and $f^{-1}(U)$ is
closed in $h^{-1}(U)$, it suffices to prove that if $\eta$ is the generic
point of $Y$ (and of $U$), we have $g(\eta)\in f^{-1}(U)$. Now, this
follows from the fact that the restriction of $g$ to $Y'$ is a $Y'$-section of
$X'\subset f^{-1}(U)$ and from the fact that $\eta\in Y'$.
\end{proof}

\paragraph{(Err\textsubscript{\textmd{IV}}, 31).}
In ({\hyperref[subsection:IV.10.4]{\textbf{IV}, 10.4}}), after Proposition {\hyperref[IV.10.4.11]{(10.4.11)}}, add:

\begin{corollary}[10.4.12]
\label{IV.10.4.12}
Let $S$ be a prescheme, $X$ an $S$-prescheme of finite presentation,
$f$ a radicial $S$-endomorphism of $X$.
\begin{enumerate}[label=\textup{(\roman*)}]
\item The endomorphism $f$ is bijective and \textup{finite (hence a \emph{universal
homeomorphism} {\hyperref[IV.2.4.5]{(2.4.5)}})}.
\item Suppose moreover that $X$ is reduced, that the local rings of
$S$ are Noetherian and universally Japanese, and that $f$ is
birational \textup{(a condition automatically satisfied if for every
maximal point $x$ of $X$, $\boldsymbol{k}(x)$ is a \emph{perfect} field)}. Then $f$ is an
$S$-automorphism.
\end{enumerate}
\end{corollary}

\begin{proof}
\textup{(i)} The question being local on $S$, we can restrict ourselves to the case
where $S=\operatorname{Spec}(A)$ is affine. There then exists a
sub-$\mathbb Z$-algebra of finite type $A_0$ of $A$ such that, on putting
$S_0=\operatorname{Spec}(A_0)$, we have
$X=X_0\times_{S_0}S$, where $X_0$ is a prescheme of finite type over $S_0$,
and that $f$ is deduced from an $S_0$-endomorphism $f_0$ of $X_0$ by the
base change $S\to S_0$
{\hyperref[IV.8.9.1]{(8.9.1)}}; we can moreover suppose
$f_0$ radicial ({\hyperref[IV.8.10.5]{8.10.5}}, (vii)). It
will suffice to prove that $f_0$ is bijective and finite, and we can therefore
suppose from now on that $A$ is a $\mathbb Z$-algebra of finite type, which
implies that $S$ and $X$ are \emph{excellent} preschemes
({\hyperref[IV.7.8.3]{(7.8.3)}} and
{\hyperref[IV.7.8.6]{(7.8.6)}}) and in particular that $X$
is a \emph{universally Japanese} prescheme
({\hyperref[IV.5.11.4]{(5.11.4)}} and
({\hyperref[IV.7.8.3]{7.8.3}}, (vi))). We can moreover
suppose $X$ \emph{reduced}; indeed,
$f_{\mathrm{red}}:X_{\mathrm{red}}\to X_{\mathrm{red}}$ is a radicial
$S_{\mathrm{red}}$-endomorphism (\textbf{I}, 5.1.6); if we prove that
$f_{\mathrm{red}}$ is finite, the same will hold for $f$: indeed,
since $f$ is radicial, it is separated
{\hyperref[IV.1.8.7.1]{(1.8.7.1)}}; since
$f_{\mathrm{red}}$ is proper, so is $f$ (\textbf{II}, 5.4.6);
but since $f$ is obviously quasi-finite, it is finite (\textbf{III}, 4.4.2).

\oldpage[IV]{354}
Let then $X'$ be the normalization of $X$ (\textbf{II}, 6.3.8); if $X_i$
($1\leq i\leq n$) are the reduced sub-preschemes of $X$ having as
underlying spaces the irreducible components of $X$, $X'$ is the
sum of the normalizations $X'_i$ of the integral preschemes $X_i$; we
shall denote by $x_i$ (resp. $x'_i$) the generic point of $X_i$
(resp. $X'_i$). The canonical morphism $p:X'\to X$ is \emph{finite} since
$X$ is universally Japanese, and we have
$p^{-1}(x_i)=\{x'_i\}$, the homomorphism $\boldsymbol{k}(x_i)\to\boldsymbol{k}(x'_i)$ being
bijective.

Let $f=(\psi,\theta)$; we already know
{\hyperref[IV.10.4.11]{(10.4.11)}} that $\psi$ is a
continuous bijection of the space $X$ onto itself; if we put
$y_i=\psi(x_i)$, $y_i$ is also a maximal point of $X$, for if $y'_i$
were a generization of $y_i$ distinct from $y_i$, $x_i$ would be a
generization $\neq x_i$ of $\psi^{-1}(y'_i)$. Since the set $M$
of maximal points of $X$ is finite, the restriction $\psi|M$ is a
\emph{bijection} of $M$ onto itself, and we can write
$\psi(x_i)=x_{\sigma(i)}$, where $\sigma$ is a permutation of
$[1,n]$. Then, for every $i$, the restriction $f|X_i$ of $f$
factors as $X_i\xrightarrow{f_i}X_{\sigma(i)}\longrightarrow X$
(\textbf{I}, 5.2.2), where the second arrow is the canonical injection. Since the
$X_i$ are integral, there exists (\textbf{II}, 6.3.9) a unique morphism
$f'_i:X'_i\to X'_{\sigma(i)}$ making the diagram
\[
\begin{array}{ccc}
X'_i&\xrightarrow{f'_i}&X'_{\sigma(i)}\\
\llap{$p_i$\,}\downarrow&&\downarrow p_{\sigma(i)}\\
X_i&\xrightarrow[f_i]{}&X_{\sigma(i)}
\end{array}
\]
commutative, and if $f'$ coincides with $f'_i$ on each of the $X'_i$, we therefore have a
commutative diagram
\phantomsection\egalabel{IV.10.4.12.1}{10.4.12.1}
\begin{equation}
\label{IV.10.4.12.1.errata}
\begin{array}{ccc}
X'&\xrightarrow{f'}&X'\\
\llap{$p$\,}\downarrow&&\downarrow p\\
X&\xrightarrow[f]{}&X
\end{array}
\tag{10.4.12.1}
\end{equation}

On the other hand, since $f$ is separated, so is
$f\circ p=p\circ f'$, hence $f'$ is \emph{separated} (\textbf{I}, 5.5.1); finally, since
$p$ is finite and $f$ bijective, it follows from
{\hyperref[IV.10.4.12.1]{(10.4.12.1)}} that $f'$ is
\emph{quasi-finite}. Since the $X'_i$ are integral and normal, it follows from
{\hyperref[IV.8.12.11]{(8.12.11)}} that $f'_i$ factors as
\[
 X'_i\xrightarrow{f''_i}X''_{\sigma(i)}
 \xrightarrow{u_{\sigma(i)}}X'_{\sigma(i)}
\]
where $X''_{\sigma(i)}$ is the integral closure of
$X'_{\sigma(i)}$ relative to the extension $\boldsymbol{k}(x'_i)$ of
$\boldsymbol{k}(x'_{\sigma(i)})$, and $f''_i$ is an open immersion. But since
$\boldsymbol{k}(x'_i)$ is a radicial extension of $\boldsymbol{k}(x'_{\sigma(i)})$,
$u_{\sigma(i)}$ is a \emph{radicial} morphism (Bourbaki,
\emph{Alg. comm.}, chap.~V, \S2, n\textsuperscript{o}~3, lemma 4);
hence $f'$ is radicial, and it then follows from
{\hyperref[IV.10.4.11]{(10.4.11)}} that $f'$ is
\emph{surjective}; we conclude (since the $u_{\sigma(i)}$ are injective)
that $f''_i$ is \emph{bijective}, hence an isomorphism; we deduce finally that
the morphism $f'$ is \emph{finite} (since this is so for $u_{\sigma(i)}$,
$X'$ being universally Japanese
{\hyperref[IV.7.8.3]{(7.8.3)}}). This being so, since
$p\circ f'=f\circ p$ is then finite and $p$ is surjective and $f$
separated, $f$ is proper (\textbf{II}, 5.4.3) and quasi-finite, hence finite
{\hyperref[IV.8.11.1]{(8.11.1)}}.

\oldpage[IV]{355}
\textup{(ii)} In view of
({\hyperref[IV.8.10.5]{8.10.5}}, (i)) applied following the
method of ({\hyperref[IV.8.1.2]{8.1.2}}, \emph{a)}), we can
restrict ourselves to the case where $S$ is a \emph{local} scheme. The preliminary reductions
made in the proof of (i) are therefore unnecessary here, and with the same
notation we see that we have a commutative diagram
{\hyperref[IV.10.4.12.1]{(10.4.12.1)}} in which $X$ is
an $S$-prescheme of finite type, $p$ is \emph{finite}, $f'$ is an \emph{isomorphism} and
the canonical homomorphism $\sh{O}_X\to p_*(\sh{O}_{X'})$ is \emph{injective}.
But we shall show in chap.~VI, in the theory of quotients of
preschemes by equivalence relations, that under these conditions $f$
is itself an isomorphism (the reader will verify that
{\hyperref[IV.10.4.12]{(10.4.12)}} will not be used in the theory of
equivalence relations of chap.~VI).
\end{proof}

\begin{corollary}[10.4.13]
\label{IV.10.4.13}
Let $k$ be a field, $X$ an algebraic prescheme over $k$. Then every
radicial $k$-endomorphism $f$ of $X$ is finite and bijective, hence a
universal homeomorphism. If moreover $X$ is reduced and $f$ birational
\textup{(for example if $k$ is of characteristic $0$)}, $f$ is a
$k$-automorphism.
\end{corollary}

\begin{remarks}[10.4.14]
\label{IV.10.4.14}
\begin{enumerate}[label=\textup{(\roman*)}]
\item In {\hyperref[IV.10.4.12]{(10.4.12)}} and
{\hyperref[IV.10.4.13]{(10.4.13)}} one cannot drop
the hypothesis that $X$ is \emph{reduced} for the second conclusion; this is
shown by the example of the $B$-endomorphism corresponding to the augmentation
$D_B(L)\to B$ of the ring $D_B(L)=A$ defined in (\textbf{0}, 18.2.3), where
$B$ is identified with a subring of $A$.

\item Let $X$ be a \emph{Noetherian} prescheme. We do not know whether every
radicial endomorphism $f$ of $X$ is necessarily surjective and finite.
If $f$ is assumed surjective and if moreover $X$ is a \emph{Japanese}
prescheme, the argument of
{\hyperref[IV.10.4.12]{(10.4.12)}} shows that $f$ is \emph{finite};
if moreover $X$ is reduced and if $f$ is birational, for example if,
for every maximal point $x\in X$, $\boldsymbol{k}(x)$ is perfect, the same
argument shows that $f$ is an automorphism.

\item One can give examples of \emph{non-surjective} radicial endomorphisms
of \emph{non-Noetherian} integral affine schemes: if $k$ is a
ring, $A=k[(T_\lambda)]_{\lambda\in L}$ a polynomial ring where $L$
is countably infinite, and $\mathfrak J$ the ideal generated in $A$ by
a subfamily $(T_\lambda)_{\lambda\in H}$, where $L\setmin H$ is infinite and
$H\neq\varnothing$, we remark that $A/\mathfrak J$ is isomorphic to
$k[(T_\lambda)]_{\lambda\in L-H}$, hence to $A$; the endomorphism of
$\operatorname{Spec}(A)$ corresponding to the composite ring
endomorphism
\[
 A\xrightarrow{p}A/\mathfrak J\xrightarrow{u}A
\]
where $p$ is the canonical projection and $u$ an isomorphism, answers the
question, its image being the closed sub-prescheme of
$\operatorname{Spec}(A)$ defined by $\mathfrak J$.
\end{enumerate}
\end{remarks}

\EGAIVErrataThirtyTwoThroughFiftyThree
